\chapter{Materiais e métodos}
\label{cap:03}
De acordo com a natureza e os objetivos do presente trabalho, foram adotadas as seguintes metodologias de pesquisa para
embasar o desenvolvimento do sistema especialista educacional proposto.

\section{Natureza da Pesquisa}
\label{sec:natureza-pesquisa}

\subsection{Bibliográfica}
Durante o desenvolvimento deste trabalho, aplica-se a metodologia bibliográfica, por meio da qual foi realizada
a análise de materiais já publicados. Para isso, fez-se necessária a avaliação e fundamentação sobre os conteúdos
consultados podendo eles ser "de livros, periódicos, documentos, textos,
mapas, fotos, manuscritos e, até mesmo, de material
disponibilizado na internet etc" \cite{fontelles2009}.

Em primeira instância, a pesquisa foi conduzida com a definição dos temas e objetivos necessários, com o intuito de identificar
os assuntos que sustentam a narrativa do trabalho, esclarecendo quais pontos são úteis e relevantes para a pesquisa. Os temas
centrais investigados foram: Problemas de Satisfação de Restrições (CSP), algoritmos de propagação de restrições,
sistemas especialistas, inteligência artificial explicável (XAI), sistemas tutores inteligentes e arquitetura de software
orientada a serviços. Além disso, foi necessário delimitar um período para a seleção das fontes, priorizando materiais mais
recentes, sem desconsiderar obras clássicas e fundamentais para o domínio, como os trabalhos de \textcite{Clancey1983},
\textcite{Mackworth1977} e \textcite{Russell2010}, dada a relevância teórica permanente de tais referências.

Após essa etapa inicial, foram selecionadas como principais bases de busca o Google Scholar e o Portal de Periódicos da CAPES,
acessado por meio do sistema CAFe (Comunidade Acadêmica Federada), em razão de sua vasta cobertura de artigos científicos,
livros, teses, dissertações e demais produções acadêmicas indexadas. Foram também realizadas buscas direcionadas em portais de
conferências e periódicos especializados em inteligência artificial, engenharia de \textit{software} e sistemas de informação.

Dado o volume de resultados retornados, realizou-se uma triagem dos materiais recuperados por meio da leitura de resumos.
Os materiais considerados relevantes foram incluídos, nos casos de dúvidas quando à permanência, procedeu-se a uma leitura
mais minuciosa. Após a seleção, foi realizada uma análise detalhada dos materiais e a síntese das informações coletadas.

Por fim, todas as pesquisas e extrações de informações foram revisadas, de modo a evitar erros e má interpretação. A revisão foi
conduzida em duas etapas, com releitura do material, da síntese e da extração de informação, resultando em ajustes quando
necessário e garantindo a qualidade e a precisão das informações utilizadas ao longo do trabalho.

\subsection{Aplicada}
A pesquisa aplicada orienta o desenvolvimento prático de compatibilidade, cujo produto é um sistema especialista educacional
para validação de compatibilidade de \textit{hardware}. Segundo \textcite{gerhardt2009}, a pesquisa aplicada tem como objetivo
gerar conhecimento para aplicação prática e dirigidos à solução de problemas específicos. Neste contexto, os fundamentos
teóricos levantados pela pesquisa bibliográfica, notadamente o formalismo de CSP, o algoritmo AC-3 e os princípios de
sistemas especialistas com \textit{explanation facility}, foram diretamente empregados na concepção e implementação do
sistema proposto.

Essa abordagem permitiu que as decisões arquiteturais e de projeto fossem fundamentadas em evidências da literatura, como a
adoção da arquitetura em três camadas \cite{Fowler2002}, o estilo REST para comunicação entre camadas \cite{Fielding2000}
e o algoritmo AC-3 como motor de inferência \cite{Mackworth1977}, garantindo o alinhamento entre os resultados práticos
obtidos e o referencial teórico consolidado.

\section{Materiais}
\label{sec:materiais}

\subsection{Equipamentos}
Para a realização do desenvolvimento, codificação e testes do sistema, foi utilizado um computador com processador Ryzen 7 5700X, 32GB de
memória RAM, placa de vídeo AMD Radeon RX 7600 e sistema operacional Windows 11. Esses recursos foram suficientes para executar o ambiente de
desenvolvimento, o servidor de aplicação, o sistema de gerenciador de banco de dados e os navegadores necessários para a validação da
interface.

\subsection{Softwares}
Para o desenvolvimento do sistema foram utilizadas ferramentas de \textit{software} selecionadas com base em sua adequação técnica ao
projeto, disponibilidade gratuita ou de código aberto e familiaridade prévia do desenvolvedor.

\section{Ferramentas Utilizadas}

\subsection{Visual Studio Code}
O Visual Studio Code (VS Code) é um editor de código-fonte desenvolvido pela Microsoft, disponível gratuitamente para sistemas
operacionais Windows, Linux e MacOS. Por meio de extensões, a ferramenta agrega funcionalidades como depuração, controle de versão,
terminal integrado e suporte a diversas linguagens de programação, aproximando-se de um ambiente de desenvolvimento integrado (IDE).

No presente trabalho, o VS Code foi utilizado como ambiente principal de desenvolvimento, tanto para codificação do \textit{backend}
e do \textit{frontend} quanto para a edição e compilação do documento LaTeX deste trabalho, por meio da extensão LaTeX Workshop.
A escolha pela ferramenta deve-se à sua leveza, à vasta disponibilidade de extensões relevantes para o ecossistema adotado
e ao conhecimento prévio do desenvolvedor.

\subsection{Node.js e NestJS}
O Node.js é um ambiente de execução JavaScript de código aberto e multiplataforma, construído sobre o motor V8 do Google Chrome,
que permite a execução de código JavaScript no lado do servidor, sendo amplamente empregado no desenvolvimento de aplicações \textit{web}
escaláveis e de alta performance. O NestJS é um \textit{framework} para desenvolvimento de aplicações server-side organizado em módulos, \textit{controllers}
e \textit{services}. Essa estrutura favorece a separação de responsabilidades e a manutenibilidade do código, alinhando-se diretamente
aos princípios da arquitetura em camadas descrito por \textcite{Fowler2002}.

No contexto do sistema proposto, o Node.js com NestJS compõe a camada de lógica de negócio, sendo o responsável por receber as requisições
da interface, executar o algoritmo AC-3 sobre o grafo de restrições de \textit{hardware} e retornar ao cliente os resultados da validação
acompanhados das justificativas educativas correspondentes.

\subsection{PostgreSQL e Docker}
O PostgreSQL é um banco de dados relacional de código aberto, com suporte a transações, multiplataforma e amplamente utilizado em aplicações
\textit{web} de diferentes escalas. A escolha por esse sistema deve-se à sua robustez, à sua conformidade com o padrão SQL e à facilidade
de integração com ecossistemas Node.js.

No sistema proposto, o PostgreSQL é utilizado para armazenar a base de conhecimento do sistema especialista: os modelos de componentes de
\textit{hardware}, suas especificações técnicas e as regras de compatibilidade que compõem o conjunto de restrições do CSP. Essas informações
são consultadas pela camada de lógica de negócio a cada execução do algoritmo AC-3.

O provisionamento do banco de dados foi realizado por meio do Docker, uma plataforma de código aberto para criação e execução de contêineres
de \textit{software}. A utilização do Docker permite que o ambiente do banco de dados seja definido de forma declarativa e reproduzível,
garantindo consistência entre os ambientes de desenvolvimento e produção sem a necessidade de instalação direta do PostgreSQL no sistema operacional.

\subsection{DBeaver}
O DBeaver é uma ferramenta de administração de banco de dados de código aberto e multiplataforma, compatível com diversos sistemas gerenciadores,
incluindo o PostgreSQL. A ferramenta oferece uma interface visual para a execução de consultas SQL, navegação entre tabelas, modelagem de dados e
exportação de resultados.

No presente trabalho, o DBeaver foi utilizado para o auxílio na administração e manipulação do banco de dados durante
o desenvolvimento, permitindo a inspeção e validação dos dados cadastrados na base de conhecimento do sistema especialista.

\subsection{Next.js}
O Next.js é um \textit{framework} de código aberto para o desenvolvimento de aplicações \textit{web} construído sobre o React, desenvolvido
e mantido pela Vercel. O \textit{framework} oferece recursos como renderização híbrida, roteamento baseado em sistemas de arquivos e
suporte nativo a APIs, facilitando o desenvolvimento de aplicações \textit{web} com uma única base de código.

No presente trabalho, o Next.js foi utilizado para o desenvolvimento da camada de apresentação do sistema, implementada como uma interface
no formato \textit{wizard} que guia o usuário, passo a passo, na seleção dos componentes de \textit{hardware}. A escolha do Next.js
deve-se à sua capacidade de integrar a interface e as chamadas à camada de lógica de negócio em um único projeto, além de atualizar
dinamicamente os componentes exibidos na tela em resposta ao resultado das validações retornadas pelo motor de inferência.

\subsection{Figma}
O Figma é uma ferramenta de \textit{design} de interface baseada em nuvem, acessada por meio de navegador \textit{web}, que permite
a criação de protótipos e \textit{layouts} de forma colaborativa em tempo real, durante o processo de criação é possível fazer uso
de diferente modelos e ter uma visibilidade geral do projeto.

No presente trabalho, o Figma foi utilizado para o desenvolvimento dos protótipos das interfaces do sistema, permitindo planejar e validar
o \textit{layout} da interface \textit{wizard} antes da sua implementação na camada de apresentação. A escolha pela ferramenta deve-se
à sua disponibilidade gratuita para uso acadêmico e ao conhecimento prévio do desenvolvedor.

\subsection{Astah Community}
O Astah Community é uma ferramenta de modelagem UML desenvolvida pela empresa japonesa Change Vision, que possibilita a criação de diagramas
de casos de uso, de classes, de sequência, de domínio e de atividades, entre outros. A ferramenta disponibiliza uma versão gratuita voltada
ao uso acadêmico.

No presente trabalho, o Astah Community foi utilizado para a elaboração dos diagramas UML, que compõe as etapas de análise
e projeto de \textit{software}, incluindo o diagrama de casos de uso, de domínio, de classes e de sequência (por enquanto).

\subsection{Teste?}

\section{Métodos de Desenvolvimento}

\subsection{Modelagem do Problema como CSP}
A etapa central do método de desenvolvimento consistiu na formalização do problema de compatibilidade de \textit{hardware}
como um Problema de Satisfação de Restrições (CSP). Conforme descrito por \textcite{Russell2010}, um CSP é caracterizado
por uma tripla $(X, D, C)$, na qual $X$ representa o conjunto de variáveis, $D$ os respectivos domínios e $C$ o conjunto
de restrições.

No sistema proposto, cada categoria de componente de \textit{hardware} foi mapeada a uma variável do CSP, cujo domínio é
composto pelos modelos comercialmente disponíveis cadastrados na base de conhecimento. As restrições foram identificadas a partir
das especificações técnicas oficiais dos fabricantes e formalizadas como restrições binárias entre pares de variáveis.

\subsection{Implementação do Algoritmo AC-3}
O motor de inferência do sistema foi implementado com base no algoritmo AC-3 (\textit{Arc Consistency Algorithm 3}),
proposto por \textcite{Mackworth1977}. A implementação seguiu a estrutura algorítmica descrita por \textcite{Russell2010},
uma fila de arcos é inicializada com todos os pares de variáveis relacionadas por restrições binárias, para cada arco retirado
da fila, o algoritmo verifica se todos os valores do domínio da variável de origem possuem suporte no domínio da variável de destino,
valores sem suporte são removidos do domínio e, a cada remoção, os arcos afetados são reinseridos na fila para propagação.

No contexto do sistema proposto, o AC-3 é invocado a cada seleção de componente realizada pelo usuário na interface. O estado
atual das seleções é enviado à camada de lógica de negócio, implementada com NestJS, que executa a propagação de restrições e
retorna ao cliente os domínios atualizados, juntamente com as justificativas técnicas referentes a cada valor removido.
Essa abordagem garante que a validação ocorra de forma incremental e em tempo real, tornando viável a execução sobre catálogos
de tamanho realista sem ocorrer explosão combinatória, o que aconteceria no caso de validação por força bruta.

\subsection{Implementação das Justificativas Educativas}
Para cada restrição violada pelo AC-3, o sistema aciona um módulo de geração de justificativas, responsável por formular,
em linguagem natural, a explicação técnica correspondente à incompatibilidade identificada. Esse módulo implementa o conceito de
\textit{explanation facility} descrito por \textcite{Clancey1983}, que distingue um sistema especialista de um mecanismo de filtragem
simples pela sua capacidade de articular a cadeia de regras e restrições que levou a uma determinada conclusão.

As justificativas foram estruturadas de forma a identificar explicitamente qual restrição foi violada, quais componentes envolvidos
e qual o princípio técnico subjacente à incompatibilidade. Por exemplo, ao detectar a seleção simultânea de um processador com
plataforma AM5 e um módulo de memória DDR4, o sistema exibe ao usuário a seguinte justificativa: \textit{"Este módulo de memória é
  incompatível: a placa-mãe compatível com o processador selecionado suporta exclusivamente DDR5, e este módulo utiliza o padrão DDR4"}.
Dessa forma, cada bloqueio é transformado em uma oportunidade de aprendizagem sobre as restrições físicas que regem a montagem de
computadores.

\subsection{Processo de Desenvolvimento de Software}
O desenvolvimento do sistema seguiu uma abordagem iterativa e incremental, na qual as funcionalidades foram implementadas e validadas
em etapas sucessivas. Em cada iteração, foram realizadas as atividades de análise de requisitos, projeto, implementação e teste, permitindo
a identificação e correção precoce de inconsistências entre os requisitos especificados e o comportamento do sistema.

A organização das tarefas foi realizada por meio de uma ferramenta de gerenciamento de projetos, com atividades distribuídas conforme seu estado
de desenvolvimento: a fazer, em desenvolvimento, em teste, concluído e adiado. Essa estrutura permitiu o acompanhamento contínuo do progresso
do projeto e a priorização das atividades ao longo do ciclo de desenvolvimento.

\subsection{Estratégia de Testes}
\label{sec:estrategia-testes}

A verificação do sistema foi conduzida em duas frentes complementares. A primeira, de
natureza funcional, destina-se a confirmar que o sistema se comporta em conformidade com
os requisitos funcionais especificados na etapa de elicitação. A segunda, de natureza
quantitativa, destina-se a avaliar empiricamente o desempenho do motor de inferência,
verificando o atendimento aos requisitos não funcionais de eficiência (RNF-01 e RNF-02)
e fornecendo evidência para a questão de pesquisa relativa à viabilidade da validação
sem recurso à avaliação exaustiva por força bruta.

\subsubsection{Validação Funcional}

A validação funcional foi conduzida por meio de casos de teste elaborados a partir dos
requisitos funcionais identificados no Capítulo~\ref{cap:05} e registrados na
ferramenta Testlink. Os casos de teste contemplam os cenários de compatibilidade e de
incompatibilidade entre componentes, a exibição das justificativas educativas associadas
a cada bloqueio e o comportamento da interface wizard nas diferentes etapas de seleção.
Para cada caso de teste foram documentados o objetivo, as pré-condições necessárias à sua
execução, os passos a serem seguidos e os resultados esperados. A execução foi realizada
de forma manual, com os resultados registrados na própria ferramenta, o que possibilita a
rastreabilidade entre os requisitos especificados e os comportamentos observados durante
a validação.

\subsubsection{Avaliação de Desempenho do Motor de Inferência}

A avaliação de desempenho tem por objetivo comprovar empiricamente que o motor de
inferência baseado no algoritmo AC-3 valida configurações de \textit{hardware} sem recorrer à
avaliação exaustiva por força bruta, cuja inviabilidade decorre da natureza NP-completa do
CSP, conforme discutido na Seção~\ref{sec:complexidade}. Para isso, o motor AC-3 foi
comparado a uma implementação de referência por força bruta, que percorre o produto
cartesiano dos domínios das cinco categorias de componentes, avaliando cada atribuição
completa candidata.

Com o intuito de assegurar a equidade da comparação, a implementação por força bruta
reutiliza o mesmo avaliador de restrições empregado pelo AC-3. Dessa forma, a diferença de
custo observada entre os dois motores é atribuível exclusivamente à estratégia de percurso
do espaço de busca, e não a variações na avaliação das regras de compatibilidade. Em
conformidade com os requisitos RNF-02 e RNF-06, a força bruta foi implementada estritamente
como instrumento de avaliação experimental, isolada em um módulo externo à aplicação, sem
exposição por meio da API REST e sem qualquer acesso pela interface do usuário.

Os cenários de avaliação foram gerados de forma sintética e controlada, tendo o tamanho do
domínio de cada categoria, denotado por $d$, como única variável independente do
experimento. Os campos categóricos, a saber, o soquete e o padrão de memória, foram gerados
de maneira determinística por rodízio (\textit{round-robin}) sobre as plataformas
suportadas, o que assegura cobertura uniforme das plataformas e preserva a correlação entre
soquete e padrão de memória fiel às plataformas reais, em conformidade com o requisito
RNF-09. Os campos quantitativos, correspondentes ao TDP e à potência, foram amostrados por
gerador pseudoaleatório inicializado com semente fixa. Essa construção isola o tamanho do
domínio como único fator sob variação e assegura o determinismo exigido pelo requisito
RNF-10, garantindo a reprodutibilidade integral do experimento.

A avaliação foi conduzida sob dois regimes, correspondentes a dois momentos distintos da
interação com o sistema. No regime irrestrito, os domínios de todas as categorias
permanecem completos, situação que representa o catálogo antes de qualquer seleção do
usuário. No regime de seleção parcial, uma das variáveis é fixada em um único valor,
situação que representa uma etapa intermediária do wizard, na qual o usuário já selecionou
um componente. Esse segundo regime reproduz o comportamento de um CSP dinâmico, no qual o
conjunto de valores admissíveis é progressivamente restringido a cada decisão do usuário
\cite{Mittal1989}.

O desempenho de cada motor foi quantificado por duas métricas. A métrica principal é o
número de avaliações de restrição executadas, escolhida por ser independente das
características da máquina de teste e, portanto, comparável entre diferentes ambientes de
execução. A métrica secundária é o tempo de execução, medido em milissegundos. Para cada
valor de $d$, cada motor foi executado um número fixo de repetições, adotando-se a mediana
como valor representativo, de modo a atenuar variações pontuais de medição.

Dada a natureza exponencial do espaço percorrido pela força bruta, estabeleceu-se um teto de
aborto, definido em termos de tempo máximo de execução e de número máximo de combinações
avaliadas. Quando o teto é atingido, a execução é interrompida e registrada como não
concluída, situação esperada para os valores mais elevados de $d$. Nesses casos, os valores
registrados correspondem ao custo acumulado até o ponto de interrupção, e não ao custo total
que a avaliação exaustiva demandaria. Os resultados obtidos por meio deste protocolo, para
ambos os regimes, são apresentados e discutidos no Capítulo~\ref{cap:resultados}.

% A metodologia especifica como os objetivos estabelecidos serão alcançados. As partes constitutivas da metodologia: a amostragem e as formas de coleta, de organização e de análise dos dados.

% Deem uma lida no link abaixo para ter uma ideia do que colocar. Conversar com o orientador para definição da metodologia.
% \url{https://blog.mettzer.com/metodologia-cientifica/amp/#O_QUE_E_METODOLOGIA_CIENTIFICA}
